\section*{Metodología estadística}

\subsection*{Descripción del conjunto de datos}

El conjunto se obtuvo del portal de datos abiertos de la empresa eléctrica estatal de Tamil Nadu. Comprende registros mensuales de enero de 2018 a diciembre de 2022 e incluye: fecha; consumo total (MWh); consumo sectorial (residencial, industrial, agrícola, comercial); demanda máxima (MW); energía suministrada (MWh); número de consumidores por sector; y variables meteorológicas (temperatura media, humedad y precipitación) del Tamil Nadu State Climate Portal.

\notebox{El resumen del artículo indica datos de abril de 2006 a marzo de 2023, mientras que la sección 3.1 los restringe a 2018 a 2022. Esta inconsistencia se retoma en el análisis de resultados.}

\subsection*{Preprocesamiento de datos}

\begin{itemize}
  \item \textbf{Imputación de valores faltantes:} combinación de interpolación lineal (para brechas cortas) y \textit{forward-fill} (para brechas largas).
  \item \textbf{Detección de atípicos:} $Z$-score (umbral $|z|>3$) y rango intercuartílico (IQR); los valores extremos se \textit{winsorizaron} en los percentiles 5 y 95.
  \item \textbf{Normalización:} Min-Max al rango $[0,1]$ para equiparar escalas.
  \item \textbf{Codificación temporal cíclica:} seno y coseno para mes y año, preservando la periodicidad:
\end{itemize}
\[ \text{month}_{\sin}=\sin\!\left(\frac{2\pi\cdot\text{month}}{12}\right),\qquad
   \text{month}_{\cos}=\cos\!\left(\frac{2\pi\cdot\text{month}}{12}\right) \]

\subsection*{Ingeniería de variables}

\begin{itemize}
  \item \textbf{Rezagos (\textit{lags}):} 1, 3 y 6 meses del consumo y de la demanda máxima.
  \item \textbf{Medias móviles (\textit{rolling}):} medias de 3 y 6 meses.
  \item \textbf{Razones sectoriales:} residencial/industrial y comercial/total.
  \item \textbf{Términos de interacción:} Temperatura $\times$ Humedad y Desviación de precipitación $\times$ Energía suministrada.
  \item \textbf{Métricas de cambio porcentual:} variación mensual del consumo y de la demanda.
\end{itemize}

\subsection*{Modelos estadísticos y de aprendizaje automático}

\noindent\textbf{Modelos basados en variables (regresión).}

\textit{Lasso} (regularización L1):
\[ \hat{\beta}=\arg\min_{\beta}\left\{\sum_{i=1}^{n}(y_i-X_i\beta)^2+\lambda\|\beta\|_1\right\} \]

\textit{Ridge} (regularización L2):
\[ \hat{\beta}=\arg\min_{\beta}\left\{\sum_{i=1}^{n}(y_i-X_i\beta)^2+\lambda\|\beta\|_2^2\right\} \]

\textit{XGBoost} (\textit{ensemble tree boosting}):
\[ \hat{y}_i=\sum_{k=1}^{K} f_k(x_i),\qquad \Omega(f)=\gamma T+\frac{1}{2}\lambda\|w\|^2 \]

\noindent\textbf{Modelos temporales.}

\textit{ARIMA$(p,d,q)$:}
\[ \phi(B)(1-B)^d\, y_t=\theta(B)\,\varepsilon_t \]
con $\phi(B)=1-\phi_1 B-\cdots-\phi_p B^p$ y $\theta(B)=1+\theta_1 B+\cdots+\theta_q B^q$, donde $B$ es el operador de retroceso.

\textit{VAR$(p)$} (para entrada multivariable):
\[ Y_t=c+\sum_{i=1}^{p} A_i\, Y_{t-i}+\varepsilon_t,\qquad Y_t\in\mathbb{R}^k,\; A_i\in\mathbb{R}^{k\times k},\; \varepsilon_t\sim\mathcal{N}(0,\Sigma) \]

\subsection*{Algoritmo híbrido propuesto}

El marco combina un componente de aprendizaje automático $M_{ML}\in\{\text{Lasso, Ridge, XGBoost}\}$ y uno de series temporales $M_{TS}\in\{\text{ARIMA, VAR}\}$. Sobre la matriz de variables enriquecida $X'$ se obtiene $\hat{y}_{ML}$; sobre la serie $y_t$ se obtiene $\hat{y}_{TS}$. La predicción final se selecciona minimizando una métrica de error $\varepsilon$ (RMSE o MAE) \citep{zhang2003}:
\[ \hat{y}_{final}=\arg\min_{\hat{y}\in\{\hat{y}_{ML},\,\hat{y}_{TS}\}}\varepsilon(y,\hat{y}) \]

\subsection*{Validación y selección de modelos}

\begin{itemize}
  \item Partición 80/20: entrenamiento 2018 a 2021, prueba 2022.
  \item \textbf{Walk-forward} (\textit{rolling-origin}) para ARIMA y VAR, preservando la integridad temporal.
  \item \textbf{Grid search + $k$-fold} ($k=5$) para Lasso, Ridge y XGBoost; el parámetro $\lambda$ se optimizó en escala logarítmica.
  \item \textbf{Estacionariedad:} test ADF (se diferenció si $p>0.05$).
  \item \textbf{Selección de rezagos:} AIC, BIC y HQIC, apoyada en ACF y PACF.
  \item \textbf{Diagnóstico de residuos:} test de Ljung-Box.
  \item \textbf{Métricas de desempeño:} MAE, RMSE, MAPE e intervalos de confianza al 95\,\%.
\end{itemize}
